% Generated draft: requires mathematical and source review.
\documentclass[11pt]{article}
\usepackage[margin=25mm]{geometry}
\usepackage{fontspec}
\setmainfont{Latin Modern Roman}
\usepackage{amsmath,amssymb,mathtools,unicode-math}
\setmathfont{Latin Modern Math}
\usepackage{enumitem,needspace,xurl,hyperref}
\setlength{\parindent}{0pt}
\setlength{\parskip}{5pt}
\setlength{\emergencystretch}{3em}
\newcommand{\N}{\mathbb{N}}
\newcommand{\Z}{\mathbb{Z}}
\newcommand{\Q}{\mathbb{Q}}
\newcommand{\R}{\mathbb{R}}
\newcommand{\C}{\mathbb{C}}
\newcommand{\F}{\mathbb{F}}
\newcommand{\A}{\mathbb{A}}
\newcommand{\PP}{\mathbb{P}}
\newcommand{\E}{\mathbb{E}}
\newcommand{\eps}{\varepsilon}
\newcommand{\dd}{\,\mathrm{d}}
\begin{document}
\section*{UnsolvedMath ID: 30001020}
\textbf{Source problem code:} catalogue entry
\subsection{Definitions and mathematical statement}
Let \((X/\Delta)\) be a smooth projective orbifold over \(\C\) of general type. In the usual divisor-pair description one writes \(\Delta=\sum_i(1-1/m_i)D_i\), with orbifold multiplicities \(m_i\), and general type means that the orbifold canonical divisor \(K_X+\Delta\) is big. The precise allowed multiplicities and smoothness convention are not specified by the extraction.

The source asks whether there is a single proper subvariety \(Y\subsetneq X\), depending on \((X,\Delta)\) and not on the map, such that every orbifold morphism
\[
f:\C\longrightarrow(X/\Delta)
\]
satisfies \(f(\C)\subseteq Y\).

Two qualifications are essential. First, the text says ``every orbifold morphism'' without restricting to nonconstant maps. If constant morphisms to ordinary points are permitted, this quantifier imposes a condition on those points as well, so it cannot simply be identified with the customary nonconstant-entire-curve formulation. Whether nonconstancy was omitted must be checked; it is not silently inserted here. Second, the orbifold-morphism convention is unstated: contact orders along \(D_i\) may be required to be at least \(m_i\) or divisible by \(m_i\), and components with infinite multiplicity require an avoidance convention. These give different classes of maps. The intended statement must select one before the common exceptional-subvariety assertion is definitive.
\subsection{Short English statement}
The proposed orbifold Green--Griffiths--Lang assertion seeks one proper algebraic subset containing all relevant entire-curve images. The extraction omits a nonconstant qualifier and leaves the orbifold contact-order convention undefined, so these are explicit source issues.
\subsection{Sources}
\begin{itemize}
\item Local UnsolvedMath registry, permanent ID \(30001020\), source code \texttt{catalogue entry}, ``Entire Curves in General Type Orbifolds.'' Registry curation: Ulam AI / UnsolvedMath Contributors, CC BY 4.0; underlying sources retain their supplied provenance.
\item \emph{Komplexe Analysis} (2009), Oberwolfach Reports, pp.\ 2165--2218, DOI \url{https://doi.org/10.4171/owr/2008/38}.
\item Demailly, supplied title \emph{Towards the Green Griffiths Lang conjecture}: \url{https://arxiv.org/abs/1412.2986}.
\item Supplied literature triage, \texttt{literature\_checked\_at}: 2026-08-21: labeled open; describes the orbifold Green--Griffiths--Lang statement as a general conjecture, with surface and other special cases known and the full smooth projective orbifold assertion open. This is supplied dated curation, not independently verified, and does not resolve the literal constant-map issue.
\end{itemize}
\end{document}
